\section{Scalar saturation and the orthogonal orbit lift}
\label{sec:tpc52-scalar-saturation}

This section separates two operators that have very different lower-frame
behaviour.  The first forgets the orbit label and adds a growing collection
of nearby scalar profiles in one Hilbert space.  The second retains that
label as an orthogonal output coordinate.  The saturation theorem below
concerns the first, unreduced Gram matrix.  It is not, by itself, a statement
about the smallest positive eigenvalue after quotienting exact relations.

\subsection{A clustered smooth curve has a small Gram eigenvalue}

Let \(\mathsf H\) be a complex Hilbert space.  For vectors
\(v_0,\ldots,v_r\in\mathsf H\), we use the Gram convention
\begin{equation}
 \operatorname{Gram}(v_0,\ldots,v_r)_{k\ell}
 :=\langle v_\ell,v_k\rangle_{\mathsf H}.
 \label{eq:tpc52-Gram-convention}
\end{equation}
Thus, for every \(c\in\C^{r+1}\),
\begin{equation}
 c^*\operatorname{Gram}(v_0,\ldots,v_r)c
 =\left\|\sum_{\ell=0}^r c_\ell v_\ell\right\|_{\mathsf H}^2.
 \label{eq:tpc52-Gram-quadratic-form}
\end{equation}

\begin{theorem}[Clustered-curve saturation]
\label{thm:tpc52-scalar-saturation}
Let \(I\subset\R\) be an interval, let \(r\geq1\), and let
\(g\in C^r(I;\mathsf H)\) satisfy \(\|g(a)\|_{\mathsf H}=1\) on
\(I\).  Fix a compact interval \(I_1\Subset I\).  Suppose that, for
some \(J\geq1\), there are \(r+1\) distinct points
\(a_0,\ldots,a_r\in I_1\) and a point \(a_*\in I_1\) such that
\begin{equation}
 \max_{0\leq\ell\leq r}|a_\ell-a_*|\leq \frac{D}{J}
 \label{eq:tpc52-cluster-width}
\end{equation}
for a fixed \(D>0\).  Then
\begin{equation}
 \lambda_{\min}\!\left(
   \operatorname{Gram}(g(a_0),\ldots,g(a_r))
 \right)
 \leq
 \frac{(r+1)D^{2r}}{(r!)^2J^{2r}}
 \sup_{a\in I_1}\|g^{(r)}(a)\|_{\mathsf H}^2.
 \label{eq:tpc52-clustered-Gram-bound}
\end{equation}
Consequently, the least eigenvalue of any larger normalized Gram matrix
containing these vectors is bounded by the right-hand side of
\eqref{eq:tpc52-clustered-Gram-bound}.
\end{theorem}

\begin{proof}
The homogeneous system
\begin{equation}
 \sum_{\ell=0}^r c_\ell(a_\ell-a_*)^q=0,
 \qquad 0\leq q<r,
 \label{eq:tpc52-vanishing-moments}
\end{equation}
has \(r\) equations and \(r+1\) unknowns.  Choose a nonzero solution
and normalize it by \(\sum_\ell|c_\ell|^2=1\).  Taylor's formula with
integral remainder gives
\begin{equation}
 g(a_\ell)
 =\sum_{q=0}^{r-1}
   \frac{g^{(q)}(a_*)}{q!}(a_\ell-a_*)^q+R_\ell,
 \qquad
 \|R_\ell\|_{\mathsf H}
 \leq \frac{|a_\ell-a_*|^r}{r!}
       \sup_{a\in I_1}\|g^{(r)}(a)\|_{\mathsf H}.
 \label{eq:tpc52-Hilbert-Taylor}
\end{equation}
The line segment between \(a_*\) and every \(a_\ell\) lies in
\(I_1\), because \(I_1\) is an interval.  The polynomial terms in
\eqref{eq:tpc52-Hilbert-Taylor} cancel by
\eqref{eq:tpc52-vanishing-moments}.  Hence Cauchy--Schwarz and
\eqref{eq:tpc52-cluster-width} give
\begin{align}
 \left\|\sum_{\ell=0}^r c_\ell g(a_\ell)\right\|_{\mathsf H}
 &\leq \sum_{\ell=0}^r|c_\ell|\,\|R_\ell\|_{\mathsf H}\notag\\
 &\leq \frac{\sqrt{r+1}D^r}{r!J^r}
       \sup_{a\in I_1}\|g^{(r)}(a)\|_{\mathsf H}.
 \label{eq:tpc52-cluster-combination}
\end{align}
The Rayleigh principle and \eqref{eq:tpc52-Gram-quadratic-form} prove
\eqref{eq:tpc52-clustered-Gram-bound}.  The assertion for a larger
Gram matrix follows either by extending \(c\) by zero or by Cauchy
interlacing.
\end{proof}

The normalization in the theorem costs no regularity on an interior
nonzero family.

\begin{corollary}[Normalization and the logarithmic orbit grid]
\label{cor:tpc52-log-grid-saturation}
Let \(f\in C^r(I;\mathsf H)\), and suppose that
\begin{equation}
 \inf_{a\in I_1}\|f(a)\|_{\mathsf H}>0
 \label{eq:tpc52-curve-nondegeneracy}
\end{equation}
on a compact interval \(I_1\Subset I\).  Then
\(g(a):=f(a)/\|f(a)\|_{\mathsf H}\) is a \(C^r\) normalized curve on
\(I_1\), and the conclusion of
\cref{thm:tpc52-scalar-saturation} holds with a constant depending only
on \(r,D\), the lower bound in
\eqref{eq:tpc52-curve-nondegeneracy}, and the \(C^r\)-bounds for \(f\).

In particular, fix \(0<z_-<z_+<\infty\), a residue modulus
\(H_0\geq1\), and put
\begin{equation}
 a_j:=\log(j/J),
 \qquad z_-J\leq j\leq z_+J,
 \qquad j\equiv a\pmod {H_0}.
 \label{eq:tpc52-log-orbit-nodes}
\end{equation}
For every fixed \(r\), any \(r+1\) consecutive nodes in this
progression form a cluster of width \(O_{r,H_0,z_-}(J^{-1})\).
Therefore their normalized scalar Gram has least eigenvalue
\begin{equation}
 \lambda_{\min}\ll_{r,H_0,z_-,f}J^{-2r}.
 \label{eq:tpc52-log-grid-Gram-decay}
\end{equation}
At the inherited scale \(J=X^{133/400+o(1)}\), the case \(r=1\)
already gives a condition-number lower bound
\(X^{133/200+o(1)}\) whenever the selected two-vector Gram is
nonsingular.  For every fixed \(r\), the corresponding lower bound is
\(X^{133r/200+o(1)}\).
\end{corollary}

\begin{proof}
The norm map is \(C^r\) away from zero, so the first assertion follows
from the chain and quotient rules.  For the second, if
\(j\asymp J\), then
\begin{equation}
 \log\frac{j+H_0}{j}
 =\log\left(1+\frac{H_0}{j}\right)
 =O_{H_0,z_-}(J^{-1}).
 \label{eq:tpc52-adjacent-log-spacing}
\end{equation}
Summing at most \(r\) such spacings verifies
\eqref{eq:tpc52-cluster-width}.  Apply
\cref{thm:tpc52-scalar-saturation}.  A normalized Gram has diagonal
entries one and hence largest eigenvalue at least one.  Dividing by
\eqref{eq:tpc52-log-grid-Gram-decay} gives the stated condition-number
bounds when the least eigenvalue is positive.
\end{proof}

\begin{remark}[Unreduced Gram versus exact quotient]
\label{rem:tpc52-unreduced-versus-quotient}
The bounds above concern the least eigenvalue of the unreduced
coefficient Gram.  If the selected vectors have an exact relation, that
least eigenvalue is zero, which is an even stronger failure of a
full-space lower frame.  Quotienting by the exact synthesis kernel
removes zero eigenvalues.  A bound for the smallest \emph{positive}
quotient eigenvalue therefore requires an additional verification that
the selected profiles have no exact relation.  The translate theorem
below supplies that verification on one important interior model; no
such verification is inferred merely from smooth dependence on the
parameter.
\end{remark}

\subsection{Exact scalar rank and metric entropy}

\begin{proposition}[Interior translates: rank and scalar-bank lower bound]
\label{prop:tpc52-translate-rank}
Let \(0\ne w\in C_c^\infty(\R)\), let
\(a_1,\ldots,a_M\in\R\) be distinct, and put
\(w_j(s):=w(s+a_j)\).  Then the \(w_j\) are linearly independent in
\(L^2(\R)\).  If an open interval \(U\) contains
\(\bigcup_{j=1}^M\supp w_j\), their restrictions are also linearly
independent in \(L^2(U)\).

Let \(V_M:=\Span\{w_1|_U,\ldots,w_M|_U\}\), so that
\(\dim_\C V_M=M\).  If \(B\subset L^2(U)\) is a linear space and a
linear map \(P:V_M\to B\) satisfies
\begin{equation}
 \|v-Pv\|_{L^2(U)}\leq\eta\|v\|_{L^2(U)}
 \qquad(v\in V_M)
 \label{eq:tpc52-uniform-relative-scalar-approximation}
\end{equation}
for some \(0\leq\eta<1\), then
\begin{equation}
 \dim_\C B\geq M.
 \label{eq:tpc52-scalar-rank-lower-bound}
\end{equation}
Moreover, if \(\mathbb B_{V_M}\) is the unit ball of \(V_M\), then
for \(0<\eta<1\) its covering number in the \(L^2(U)\) metric obeys
\begin{equation}
 \mathcal N(\mathbb B_{V_M},\eta)\geq\eta^{-2M}.
 \label{eq:tpc52-translate-entropy}
\end{equation}
\end{proposition}

\begin{proof}
If \(\sum_jc_jw(s+a_j)=0\) on \(\R\), Fourier transformation gives
\begin{equation}
 \widehat w(\xi)\sum_{j=1}^M c_j\e^{ia_j\xi}=0.
 \label{eq:tpc52-translate-Fourier-relation}
\end{equation}
The entire function \(\widehat w\) is not identically zero, and hence
is nonzero on some real open interval.  The exponential polynomial in
\eqref{eq:tpc52-translate-Fourier-relation} vanishes on that interval,
so it vanishes identically.  Its first \(M\) derivatives at one point
form a Vandermonde system in the distinct \(a_j\), forcing every
\(c_j=0\).  If the restrictions have a relation on \(U\), the linear
combination is supported in \(U\) and vanishes there, hence vanishes on
all of \(\R\).  This proves the rank assertion.

If \(Pv=0\), \eqref{eq:tpc52-uniform-relative-scalar-approximation}
gives \(\|v\|\leq\eta\|v\|\), so \(v=0\).  Thus \(P\) is injective
on the \(M\)-dimensional space \(V_M\), proving
\eqref{eq:tpc52-scalar-rank-lower-bound}.  Finally regard \(V_M\) as a
real Euclidean space of dimension \(2M\).  Orthogonally projecting the
centres of any radius-\(\eta\) cover onto \(V_M\) does not increase
distances.  Comparing \(2M\)-dimensional volumes of the unit ball and
radius-\(\eta\) balls yields
\eqref{eq:tpc52-translate-entropy}.
\end{proof}

Thus, whenever an interior scalar translate family has
\(M\asymp J\), every scalar bank satisfying a uniform relative-output
estimate with error below one has rank \(\gg J\).  At
\(J=X^{133/400+o(1)}\), a scalar bank of polylogarithmic rank cannot do
this.  This rank statement is conditional on the displayed interior
translate model; it is not being promoted to a quotient-rank theorem
for every restricted product profile.

\subsection{Retaining the orbit coordinate}

The physical coefficient space uses a different synthesis.

\begin{theorem}[Exact orthogonal-lift identity]
\label{thm:tpc52-orthogonal-lift}
Let \((\Omega,\sigma)\) be a measure space, let \(\mathsf K\) be a
complex Hilbert space, and let \(S\) be a finite set.  For
\(f_j\in L^2(\sigma;\mathsf K)\), define
\begin{equation}
 \mathsf T^\perp c(x)
 :=\bigoplus_{j\in S}c_jf_j(x)
 \in\ell^2(S;\mathsf K),
 \qquad c=(c_j)_{j\in S}\in\C^S.
 \label{eq:tpc52-orthogonal-lift-operator}
\end{equation}
Then
\begin{equation}
 \|\mathsf T^\perp c\|_{L^2(\sigma;\ell^2(S;\mathsf K))}^2
 =\sum_{j\in S}|c_j|^2
   \|f_j\|_{L^2(\sigma;\mathsf K)}^2.
 \label{eq:tpc52-orthogonal-lift-identity}
\end{equation}
Consequently its coefficient Gram is the diagonal matrix
\begin{equation}
 \diag\left(\|f_j\|_{L^2(\sigma;\mathsf K)}^2:j\in S\right).
 \label{eq:tpc52-orthogonal-lift-Gram}
\end{equation}
After discarding the zero vectors and normalizing each nonzero
\(f_j\), the Gram is exactly the identity.  All statements remain true
when \(\sigma\) is a finite weighted sampling measure.
\end{theorem}

\begin{proof}
For almost every \(x\), orthogonality of the coordinate copies of
\(\mathsf K\) gives
\begin{equation}
 \left\|\bigoplus_{j\in S}c_jf_j(x)\right\|^2
 =\sum_{j\in S}|c_j|^2\|f_j(x)\|_{\mathsf K}^2.
 \label{eq:tpc52-pointwise-orthogonal-lift}
\end{equation}
Integrate the finite nonnegative sum.  The Gram and normalization
statements follow immediately.
\end{proof}

The theorem does not reduce the number of orbit coordinates.  It says
that a row-log frequency bank may have Hilbert-valued coefficients
carrying all those coordinates without creating the scalar cross terms
from \cref{thm:tpc52-scalar-saturation}.

\begin{proposition}[Sparse deletion through leverage]
\label{prop:tpc52-mask-leverage}
Let \(\cR\) be a finite row set of size \(N\), let \(w_\gamma>0\), and
let \(M(\alpha,\gamma)\in\{0,1\}\) satisfy
\begin{equation}
 \#\{\gamma\in\cR:M(\alpha,\gamma)=0\}\leq\eps N
 \label{eq:tpc52-leverage-mask-degree}
\end{equation}
for every source label \(\alpha\).  For a finite orbit set \(S\), let
\(f_{\alpha,j}:\cR\to\mathsf K\), and suppose
\begin{equation}
 E_{\alpha,j}
 :=\sum_{\gamma\in\cR}w_\gamma
        \|f_{\alpha,j}(\gamma)\|_{\mathsf K}^2>0.
 \label{eq:tpc52-component-row-energy}
\end{equation}
Define the leverage
\begin{equation}
 \mathfrak L_{\alpha,j}
 :=\frac{N\max_{\gamma\in\cR}
       w_\gamma\|f_{\alpha,j}(\gamma)\|_{\mathsf K}^2}
       {E_{\alpha,j}},
 \qquad
 \mathfrak L_*:=\sup_{\alpha,j}\mathfrak L_{\alpha,j}.
 \label{eq:tpc52-mask-leverage}
\end{equation}
Then, for every \(c\in\C^S\),
\begin{align}
 &\sum_{j\in S}|c_j|^2
   \sum_{\gamma\in\cR}M(\alpha,\gamma)w_\gamma
       \|f_{\alpha,j}(\gamma)\|_{\mathsf K}^2
 \notag\\
 &\hspace{20mm}\geq
 (1-\eps\mathfrak L_*)
 \sum_{j\in S}|c_j|^2E_{\alpha,j}.
 \label{eq:tpc52-masked-orthogonal-lower-bound}
\end{align}
The reverse inequality with constant one is automatic.  If
\(\eps\mathfrak L_*\geq1\), the displayed lower bound is interpreted
as the trivial nonnegative bound.
\end{proposition}

\begin{proof}
For each \((\alpha,j)\), the deleted energy is at most
\begin{align}
 \#\{\gamma:M(\alpha,\gamma)=0\}
 \max_\gamma w_\gamma\|f_{\alpha,j}(\gamma)\|^2
 \leq\eps\mathfrak L_{\alpha,j}E_{\alpha,j}.
 \label{eq:tpc52-leverage-deletion-proof}
\end{align}
Subtract this estimate, multiply by \(|c_j|^2\), and sum in \(j\).
Deleting nonnegative terms proves the upper bound.
\end{proof}

\begin{remark}[Scope of the leverage condition]
\label{rem:tpc52-leverage-scope}
The exact lift \eqref{eq:tpc52-orthogonal-lift-identity} is purely
Hilbert-space algebra.  A bound \(\mathfrak L_*=O(1)\) is a separate
sampling assertion.  Later sections verify an aggregate version on
uniformly active prime--squarefree row blocks.  Neither the identity nor
the leverage estimate is a statement about coherent prime summation,
sieve parity, or a fixed-shift prime-pair lower bound.
\end{remark}
